\section{09.11.2025 -- Wortschatz}

\begin{vocabtable}
\deword{überhaupt}{at all, actually; synonym: völlig}{Das habe ich überhaupt nicht verstanden.}
\deword{das Jobangebot}{job offer; synonym: Stellenangebot}{Ich habe ein interessantes Jobangebot bekommen.}
\deword{der Bereich}{field, area, sector; synonym: Gebiet}{Ich arbeite im Bereich Softwareentwicklung.}
\deword{eigentlich}{actually, really, by the way; synonym: tatsächlich}{Eigentlich wollte ich heute lernen.}
\deword{loslegen}{to get started; synonym: anfangen}{Wir können jetzt loslegen.}
\deword{sich bewegen}{to move oneself; synonym: sich fortbewegen}{Ich bewege mich gerne an der frischen Luft.}
\deword{der Ausgleich}{balance, compensation; synonym: Balance}{Sport ist für mich ein guter Ausgleich.}
\deword{bestimmt / bestimmtes}{specific, particular, certain; synonym: gewiss}{Ich suche ein bestimmtes Buch.}
\deword{berichten über + akk}{to report, to tell; synonym: mitteilen}{Er berichtet über seine Erfahrung.}
\deword{ausdrücken}{to express, to convey; synonym: mitteilen}{Ich möchte meine Meinung klar ausdrücken.}
\deword{die Absichten}{intentions, aims; synonym: Ziele}{Seine Absichten waren unklar.}
\deword{absolvieren + akk}{to complete, e.g. a course or internship; synonym: beenden}{Ich absolviere ein Praktikum.}
\deword{fassen + akk}{to grasp, to comprehend; synonym: begreifen}{Ich kann es kaum fassen.}
\deword{bilden + akk}{to form, create; synonym: formen}{Diese Wörter bilden einen Satz.}
\deword{ergänzen}{to complete, to add; synonym: hinzufügen}{Ich möchte noch etwas ergänzen.}
\deword{die Regeln}{rules, regulations; synonym: Vorschriften}{Man muss die Regeln beachten.}
\deword{gelten für + akk}{to apply, to be valid; synonym: zutreffen}{Diese Regel gilt für alle.}
\deword{drücken + akk}{to press; also to express}{Er drückt den Knopf.}
\deword{passend}{suitable, appropriate; synonym: geeignet}{Das ist eine passende Formulierung.}
\end{vocabtable}
